\renewcommand*{\authorOfNextLines}{Helmut Quinz}
\begin {longtable} [ht] {p{40mm}p{110mm}}
\stretchTableRows{1.3}
	Datei 				 		& Inhalt\\ \hline
	\endhead
	\hline
	\endfoot
	\endlastfoot

	\texttt{0Präambel}			& Bindet die weiteren zur Gestaltung des Dokuments erforderlichen Dateien ein.\\
	\texttt{CompilerOptions}	& Legt Vorgaben für das Generieren des Dokuemnts fest\\
	\texttt{DefColors}			& Definition einiger nicht in eingebundenen Paketen enthaltener Farben. 
								  Eine Anpassung durch den Anwender ist unproblematisch.\\
	\texttt{DefMakrosAllgemein}	& Definition im Dokument verwendeter selbst erstellter Makros.
								  Eine Anpassung sollte nur von erfahrenen Anwendern vorgenommen werden.\\
	\texttt{DefPackages}		& Einbindung von Makro-Paketen.
								  Eine Anpassung sollte nur von erfahrenen Anwendern vorgenommen werden.\\
	\texttt{DefStyleDA}			& Festlegung der optischen Gestaltung.
								  Eine Anpassung sollte nur von erfahrenen Anwendern vorgenommen werden.\\
	\texttt{DoAbstract}			& Erstellen der Kurzbeschreibung.
								  Eine Anpassung sollte nur von erfahrenen Anwendern vorgenommen werden.\\
	\texttt{DoAbstractEnglish}	& Aufbau der Kurzbeschreibung in Englisch
								  Eine Anpassung sollte nur von erfahrenen Anwendern vorgenommen werden.\\
	\texttt{DoAbstractGerman}	& Aufbau der Kurzbeschreibung in Deutsch
								  Eine Anpassung sollte nur von erfahrenen Anwendern vorgenommen werden.\\
	\texttt{DoTitle}			& Erstellen der Titelseite
								  Eine Anpassung sollte nur von erfahrenen Anwendern vorgenommen werden.\\
	\texttt{TitleHelpers}		& Makros die beim Erstellen der Titelseite unterstützen.
								  Eine Anpassung sollte nicht erforderlich sein.\\
\caption {Dateien im Verzeichnis \texttt{00\_DefTex}}
\end {longtable}
\FloatBarrier
\subsubsection{Hinweis:}
Die Bilddateien in diesem Ordner werden beim Übersetzen des Dokuments eingebunden.